\documentclass[12pt,a4paper,psamsfonts]{amsart}
\usepackage{amssymb,mathtools}
\usepackage{calc}
\setlength{\topmargin}{0cm}
\setlength{\oddsidemargin}{0cm}
\setlength{\evensidemargin}{0cm}
\setlength{\marginparwidth}{0cm}
\setlength{\marginparsep}{0cm}
\setlength{\textheight}{\paperheight - 2in - 35pt}
\setlength{\textwidth}{\paperwidth - 2in}
\setlength{\headheight}{12.5pt}
\setlength{\headsep}{25pt}
\setlength{\footskip}{30pt}
\renewcommand{\baselinestretch}{1.3}
\pagestyle{headings}
\raggedbottom
\setlength{\emergencystretch}{2em}
\usepackage[colorlinks=true,urlcolor=blue,linkcolor=blue,citecolor=blue]{hyperref}
\newtheorem*{theorem}{Theorem}
\theoremstyle{definition}
\newtheorem*{setup}{Setup}
\theoremstyle{plain}
\newtheorem{lemma}{Lemma}[section]
\newtheorem{proposition}[lemma]{Proposition}
\newtheorem{corollary}[lemma]{Corollary}
\newcommand{\Q}{\mathbb Q}
\newcommand{\C}{\mathbb C}
\newcommand{\OO}{\mathcal O}
\newcommand{\mult}{\operatorname{mult}}
\newcommand{\nef}{\operatorname{n}}
\title{Nef Abundance: Framework -- Failure of Jet Separation}
\date{Temporary draft -- October 11, 2026}
\begin{document}
\maketitle

\section{Failure of Jet Separation}

Let $X$ be a smooth projective threefold over $\C$ such that
$K_X\sim0$ and $q(X)=0$. Let $D$ be a nef Cartier divisor on $X$
such that $\nef(X,D)=3$ and $D^3=0$.
Assume in addition that $\kappa(D)=-\infty$.

\begin{setup}
\leavevmode
\begin{itemize}
\setlength{\itemsep}{2pt}
\setlength{\parsep}{0pt}
\item Fix $\epsilon\in\Q_{>0}$ sufficiently small such that
\[
 343\epsilon^3<\tfrac14,\qquad40\epsilon<\tfrac38.
\]
\item Fix an ample Cartier divisor $A$.
For integers $n\ge1$, put
\[
 a_n=((D+A/n)^3)^{1/3},\qquad
 r_n=\max\{1,\lfloor\epsilon/a_n+1/2\rfloor\},\qquad
 P_n=r_n(D+A/n).
\]
Then $P_n$ is ample and $\mathbb{Q}$-Cartier, $r_n\to\infty$, and
$P_n^3\to\epsilon^3$. Indeed, $a_n\to0$ and
$|r_na_n-\epsilon|\le a_n/2$ for all sufficiently large $n$.
\item Let $p_i:X\times X\to X$ be the projections, for $i=1,2$.
Put
\[
 Z=\mathbb P_{X\times X}
       \bigl(p_1^*\OO_X(D)\oplus p_2^*\OO_X(D)\bigr).
\]
Let $\pi:Z\to X\times X$ be the projection.
Let $S_1,S_2\subset Z$ be the two summand sections and set
\[
 U=Z\setminus(S_1\cup S_2).
\]
\item Set $\xi=c_1(\OO_Z(1))$. For a positive integer $s$, put
\[
 R_{n,s}=(p_1\circ\pi)^*P_n+(p_2\circ\pi)^*P_n+s\xi.
\]
Whenever $|kR_{n,s}|$ is used, $k>0$ is an integer with
$kR_{n,s}$ Cartier, equivalently $kP_n$ Cartier.
\end{itemize}

\end{setup}

\medskip
\noindent\textbf{Jet separation.}
Let $W$ be a smooth variety, $L$ a Cartier divisor, and $z\in W$.
For an integer $j\ge0$, $|L|$ separates $j$-jets at $z$ if
\[
 \forall u\in\OO_{W,z}(L)\quad
 \exists\sigma\in H^0(W,\OO_W(L))\quad
 \sigma_z-u\in\mathfrak m_z^{j+1}\OO_{W,z}(L).
\]
Here $\mathfrak m_z\subset\OO_{W,z}$ is the maximal ideal.

\setcounter{lemma}{1}
\begin{lemma}\label{lem:full-nojets}
For every positive integer $s$ there exists
an integer $n_0=n_0(\epsilon,s)\ge1$ such that, for every integer
$n\ge n_0$, $|kR_{n,s}|$ does not separate $8k$-jets at any $z\in U$,
for any positive integer $k$ with $kR_{n,s}$ Cartier.
\end{lemma}

\medskip
\noindent\textbf{Draft status.}
The proof below is extracted from the existing threefold proof draft.
Its source-dependent inputs and arguments have not been independently
verified. This note isolates the indicated framework lemma for discussion
and checking.

\section{Proof of Lemma~\ref{lem:full-nojets}}

\subsection{The scalar input}\label{sec:geometric-inputs}

\begin{lemma}\label{lem:full-scalar}
Let $X$ be a smooth projective threefold over $\C$ with
$K_X\sim0$ and $q(X)=0$. Let $D$ be a nef Cartier divisor with
$\nef(X,D)=3$, $D^3=0$, and $\kappa(D)=-\infty$.
Fix an ample Cartier divisor $A$ and $\epsilon\in\Q_{>0}$ such that
$343\epsilon^3<1/4$ and $40\epsilon<3/8$.
For $n\ge1$, let $r_n$ be the nearest positive integer to
$\epsilon/((D+A/n)^3)^{1/3}$ and set $P_n=r_n(D+A/n)$.
For every integer $s>0$ and all $n\gg0$, let $\beta:\widetilde X\to X$
be the blowup at a very general point $x$, with exceptional divisor $E$.
There are a birational morphism $f:V\to\widetilde X$ from a smooth
projective threefold and ample Cartier divisors $H_1,H_2$ on $V$ such that
\[
 f^*\beta^*(P_n+sD)\cdot H_1\cdot H_2
 <40\epsilon\,f^*E\cdot H_1\cdot H_2.
\]
\end{lemma}

Write $\nu(D)=\max\{j:D^j\cdot A^{3-j}>0\}\in\{1,2\}$.
We obtain the curve bound from a multiplicity estimate and duality.
All auxiliary varieties in this section are over $\C$.

\subsubsection{Normalization and numerical inputs}\label{sec:lower-dimensional-inputs}

\begin{lemma}\label{lem:hd-normalization}
For every $\epsilon>0$ and
as $n\to\infty$ through the positive integers, there are
positive integers $r_n\to\infty$ such that
\[
 P_n=r_n(D+A/n),\qquad P_n^3\longrightarrow\epsilon^3.
\]
The divisor $P_n$ is ample and $\mathbb{Q}$-Cartier; its Cartier multiples
are chosen after this normalization.
\end{lemma}
\begin{proof}
Put $a_n=((D+A/n)^3)^{1/3}$. Then $a_n>0$ and $a_n\to0$.
Choose $r_n$ to be the nearest positive integer to $\epsilon/a_n$.
For sufficiently large $n$,
$|r_na_n-\epsilon|\le a_n/2$, which proves the assertions.
\end{proof}

\begin{lemma}\label{lem:lower-dimensional}
Let $(Y,B)$ be an auxiliary projective klt pair of dimension
$f\in\{1,2\}$, with $T=K_Y+B$ pseudo-effective, and let $M$ be
nef and $\mathbb{Q}$-Cartier. Then $T+M$ has an effective rational
numerical representative. Moreover, if for some $a>0$
\[
 (T+aM)\cdot C>0
\]
for every integral curve through a sufficiently very general point,
then $T+cM$ is big for every rational $c>0$.
\end{lemma}
\begin{proof}
Numerical nonvanishing on surfaces is
\cite[Corollary~C]{LP}; on curves it follows from the degree criterion. For the second assertion, the curve case is immediate.
In dimension two run the ordinary $(K_Y+B)$-minimal model program,
after a small $\Q$-factorialization. Its output is a morphism
$g:Y\to Y_0$ with $U=K_{Y_0}+B_0$ nef. The pushforward
$M'=g_*M$ is nef. Negativity gives
\[
 T=g^*U+E_T,\qquad M=g^*M'-E_M,
\]
where $E_T,E_M\ge0$ are exceptional; $E_T$ has positive
coefficient on every contracted divisor. Thus, for small rational
$t>0$, $E_T-tE_M\ge0$.

By \cite[Corollary~C]{LP}, $U+tM'$ is numerically equivalent
to a semiample rational divisor. If it were not big, the resulting
morphism would have positive-dimensional fibres. A general fibre,
or a general curve in such a fibre when the image is a point,
avoids the finitely many points $g(\operatorname{Exc}g)$.
Since $U$ and $M'$ are nef, this curve has both $U$- and $M'$-degree
zero. Its inverse image contradicts the stated full positivity of
$T+aM$. Consequently $U+tM'$, and hence
$T+tM=g^*(U+tM')+E_T-tE_M$, is big.
For $c>0$, choose $0<\lambda<\min\{1,c/t\}$ and write
\[
 T+cM=\lambda(T+tM)+(1-\lambda)T+(c-\lambda t)M.
\]
The last two summands are pseudo-effective, so the sum is big.
This proves exactly the dimension-one and dimension-two inputs
used in \cite[Sections~3--5]{NS}, without assuming threefold
nef nonvanishing.
\end{proof}

\subsubsection{Moving base components}

\begin{lemma}\label{lem:hd-base-degree}
Let $W$ be smooth projective of dimension $d$, let $H$ be a nef
Cartier divisor on $W$, and let
$\Lambda\subseteq|H|$ be a nonempty linear system
with base ideal $\mathfrak b$. Let $V$ be an integral component of
its base locus, of codimension $c\ge1$. If
\[
 \mathfrak b\OO_{W,\eta_V}\subseteq
 \mathfrak m_{\eta_V}^{h}\qquad(h\in\mathbb Z_{>0}),
\]
then
\[
 h^c\,H^{d-c}\cdot V\le H^d.
\]
In particular, if $x$ is an isolated base point and every member
has multiplicity at least a real number $r\ge0$ at $x$, then
$r^d\le H^d$.
\end{lemma}
\begin{proof}
Intersect successively with general members of $\Lambda$.
Before each next cut, remove the components already contained in
the base locus. Each remaining component is cut properly, and the
intersection cycle is effective. Its $H$-degree equals the $H$-degree
of the preceding cycle; discarding effective components can only
decrease this degree, since $H$ is nef. Starting with $[W]$, the
$H$-degree of the codimension-$c$ cycle is therefore at most $H^d$.

At $\eta_V$, the localized base ideal is primary for the maximal
ideal. No discarded component of smaller codimension contains
$\eta_V$, because $V$ is an irreducible component of the base locus.
The first $c$ general equations $s_1,\ldots,s_c$ form a parameter
ideal $I$ in the regular local ring $R=\OO_{W,\eta_V}$.
The coefficient of $V$ in the intersection cycle is
$\operatorname{length}(R/I)$. Since $I\subseteq\mathfrak m^h$,
monotonicity of Hilbert--Samuel multiplicity gives
\[
 \operatorname{length}(R/I)=e(I,R)
 \ge e(\mathfrak m^h,R)=h^c.
\]
Here the first equality uses that a parameter sequence in $R$ is
regular, and the last equality uses regularity of $R$.
This proves the inequality. For an isolated point take $c=d$ and
$h=\max\{1,\lceil r\rceil\}$.
\end{proof}

\begin{lemma}\label{lem:finite-subadjunction-model}
Let $W$ be smooth projective, $\Theta\ge0$ rational, and $V\subset W$
integral of positive dimension. Suppose $(W,\Theta)$ is lc at
$\eta_V$ and has an lc place centred on $V$.
For $L=(K_W+\Theta)|_{V^\nu}$ there exist a finite surjective
morphism $\eta:V'\to V^\nu$, with $V'$ normal, and rational
divisors $B_{V'}\ge0$ and $M_{V'}$, the trace of a b-nef divisor,
such that
\[
 \eta^*L\sim_\Q K_{V'}+B_{V'}+M_{V'}.
\]
\end{lemma}
\begin{proof}
We use the construction in \cite[Lemma~7.1]{LA}.
The dlt blowup \cite[Theorem~2.10]{FH} gives
\[
 p^*(K_W+\Theta)=K_Q+B_Q+E_Q,
 \qquad (Q,B_Q)\text{ dlt},\quad E_Q\ge0,
\]
where $B_Q$ is the truncated boundary and the effective surplus
$E_Q$ is absent near $\eta_V$. Choose a dlt stratum $S$ minimal
among those dominating $V$. Adjunction, followed by adding
$E_Q|_S$, gives an effective rational divisor $B_S$ and
\[
 K_S+B_S\sim_\Q(p|_S)^*((K_W+\Theta)|_V).
\]
The stratum is not contained in $E_Q$. Minimality makes
$(S,B_S)$ klt over $\eta_V$, although it need not be globally lc.

Factor through the normalization and the connected-fibre base:
$S\xrightarrow{q}V'\xrightarrow{\eta}V^\nu\to V$, with $\eta$
finite. The discrepancy-round-up rank condition for $q$ is one:
over the generic point the effective klt boundary contributes
zero round-up on strict divisors and allows poles only along
exceptional divisors. Its pushforward to normal $S$ is therefore
$\OO_S$, and connectedness gives rank one. Thus this is a
klt-trivial fibration in \cite[Definition~3.1]{FG}. Its formula is
\[
 \eta^*((K_W+\Theta)|_{V^\nu})
 \sim_\Q K_{V'}+B_{V'}+M_{V'},
\]
with b-nef moduli divisor $M$ by \cite[Theorem~3.5]{FG}.
The discriminant on this initial base is effective: over a prime
divisor, an inverse-image component of multiplicity $m\ge1$ and
boundary coefficient $b\ge0$ bounds its lc threshold by
$(1-b)/m\le1$. A negative threshold still gives a positive
discriminant coefficient. Only generic sub-klt singularities are
required for this use of the formula.
\end{proof}

\begin{lemma}\label{lem:generic-subadjunction}
Under the hypotheses of Lemma~\ref{lem:finite-subadjunction-model},
let $f=\dim V$, let $P_V$ be nef and $\mathbb{Q}$-Cartier, and let
$\rho:Y\to V$ be a resolution. Then
\[
 K_Y\cdot (\rho^*P_V)^{f-1}\le(K_W+\Theta)|_V \cdot P_V^{f-1}.
\]
\end{lemma}
\begin{proof}
Apply Lemma~\ref{lem:finite-subadjunction-model} and intersect with
$f-1$ copies of the pulled-back $P_V$. The effective discriminant
and the moduli trace have nonnegative intersection; for the latter,
push down from a higher base where the moduli divisor is nef.
For compatible canonical Weil cycles,
\[
 \eta_*K_{V'}=(\deg\eta)K_{V^\nu}+\eta_*R_\eta,
 \qquad R_\eta\ge0.
\]
Discard ramification and divide by $\deg\eta$. Pushing down $K_Y$
to $V^\nu$ changes it only by exceptional cycles, whose testing
intersection is zero. No $\mathbb{Q}$-Cartier assumption on
$K_{V^\nu}$ is needed.
\end{proof}

\begin{lemma}\label{lem:canonical-centre}
Let $W$ be smooth projective of any dimension, and let $P_W$
be nef and $\mathbb{Q}$-Cartier. Suppose a positive-dimensional integral base component $V$
of a subseries of $|kP_W|$ is isolated at its generic point,
of codimension $a$, and its base ideal there is contained in
$\mathfrak m^h$. Then on a resolution $\rho:Y\to V$,
\[
 K_Y\cdot (\rho^*P_W)^{\dim V-1}
 \le (K_W+cP_W)|_V \cdot P_W^{\dim V-1},
 \qquad 0<c\le ak/h.
\]
\end{lemma}
\begin{proof}
The generic log canonical threshold of the base ideal satisfies
$\lambda\le a/h$, as seen by blowing up the smooth generic
centre. On a log resolution, average sufficiently many general
members of the subseries with total coefficient $\lambda$.
This gives an actual effective boundary
$\Theta\sim_{\Q}k\lambda P_W$ which is lc at the generic
point of $V$ and has an lc place centred on $V$.
Lemma~\ref{lem:generic-subadjunction} gives the inequality with
$c=k\lambda$. The construction uses actual general members of
the given subseries, and requires lc only at $\eta_V$.
Both the threshold bound and subadjunction are dimension-independent.
\end{proof}

\begin{lemma}\label{lem:hd-dense-multiplicity}
Let $H$ be a Cartier divisor on $X$, let
$\Lambda\subseteq|H|$ be a fixed nonempty linear system, and let
$r\ge0$ be real. If a Zariski dense set of points $x$ admits
$E_x\in\Lambda$ with $\mult_x(E_x)>r$, then every point does.
\end{lemma}
\begin{proof}
Use local parameters $u,v,w$ and homogeneous parameters
$[t_0:\cdots:t_m]$ on the fixed linear system. Write the universal
equation as $f=\sum_i t_i f_i$, and denote its coefficient functions
by $c_{i,\alpha}(x)$, where $\alpha\in\mathbb Z_{\ge0}^3$.
The incidence $Z$ is defined locally by
\[
 \sum_i t_i c_{i,\alpha}(x)=0\qquad(|\alpha|\le\lfloor r\rfloor).
\]
These coefficients are the finite local jet coefficients. Thus
$Z$ is closed, and its proper projection has closed image.
An image containing a Zariski dense subset equals $X$.

\end{proof}

\begin{lemma}\label{lem:jet-normal-containment}
Let $W$ be smooth projective, $L$ a line bundle, and $z:T\to W$
a morphism from a smooth integral variety. Fix integers
$m_1>m_0\ge0$ and $1\le h\le m_1-m_0+1$.
Suppose the spaces
\[
 \mathcal H_i(t)=H^0(W,L\otimes\mathfrak m_{z(t)}^{m_i})
 \quad(i=0,1)
\]
form vector bundles, and a reduced family $\Gamma\subset T\times W$
dominates $W$, with each general fibre $\Gamma_t$ a component of
both base loci. Then the higher base ideal satisfies
$\mathfrak b_{1,t}\subset I_{\Gamma_t}^{h}$ at a general incidence point.
\end{lemma}
\begin{proof}
A frame of the higher kernel has the form
$s(t,w)=\sum_\alpha a_\alpha(t)s_\alpha(w)$ for fixed global sections.
A parameter derivative of order $j<h$ loses at most $j$ orders
at the moving mark, so belongs to the full lower kernel and
vanishes on $\Gamma$.
At a general smooth point, dominance of $\Gamma\to W$ makes
parameter directions span the normal space. Choose transverse
parameters $t'$ and write $\Gamma$ locally as
$u=t'-F(t'',w)=0$. Differentiation in $t'$ with $t'',w$ fixed
is differentiation in $u$. Thus all terms of degree less than $h$
in the normal parameters vanish, giving $s\in I_\Gamma^h$.
Characteristic zero and faithful flatness of the analytic local
ring give the algebraic containment. Restricting the finite frame
to a general parameter fibre proves the assertion, as in
\cite[Lemma~7.3]{LA}.
\end{proof}

\begin{lemma}\label{lem:restricted-components}
Let $g:W\to B$ be a morphism from a smooth projective variety,
and let an integral component $V$ of the base locus of a linear
system have base ideal $\mathfrak b\subset I_V^h$.
Let $F$ be the reduced component through a general point of a
fibre of $V\to g(V)$. Assume the ambient fibre $W_b$ is smooth
and integral and $F\subsetneq W_b$. Then the restricted system
is nonzero, $F$ is a base component, and
\[
 \mathfrak b|_{W_b}\subset I_F^h
\]
near $\eta_F$.
\end{lemma}
\begin{proof}
Remove the other base components and work where $V\to g(V)$ is
smooth over its actual image. Its scheme fibre is reduced there,
hence $I_V\OO_{W_b}=I_F$ and $I_V^h\OO_{W_b}=I_F^h$.
The local base support is exactly $F$.
Since $F$ is proper, not all restricted sections vanish identically
on $W_b$. The resulting base ideal is primary at $\eta_F$.
This is the actual-image restriction in \cite[Lemma~7.4]{LA}.
\end{proof}

\begin{lemma}\label{lem:full-kernels}
Let $W$ be a smooth projective variety of dimension $d$, and let
$P_W$ be ample and $\mathbb{Q}$-Cartier. Fix $kP_W$ Cartier and levels
$\tau_0<\cdots<\tau_d$ with common gap $\delta>0$ and
$k\delta\ge4$. Suppose that, at a very general marked point $z$,
all full spaces
\[
 \mathcal H_i(z)=H^0(W,\OO_W(kP_W)\otimes\mathfrak m_z^{\lceil k\tau_i\rceil})
\]
are nonzero, and the lowest base locus has a positive-dimensional
component through $z$. Then there is an algebraic covering family
of proper integral components $V_t$, common to two adjacent base
loci, with $1\le f=\dim V_t<d$. At a general incidence point the
entire higher-kernel base ideal satisfies
\[
 \mathfrak b_t\subset I_{V_t}^h,\qquad
 h=\lceil k\delta/2\rceil,\qquad k/h\le2/\delta.
\]
Consequently, on a resolution $V_t^*$,
\begin{equation}\label{eq:full-kernel-bounds}
 \begin{split}
 0<P_W^f\cdot V_t&\le(2/\delta)^{d-f}P_W^d,\\
 K_{V_t^*}\cdot P_W^{f-1}&\le(K_W+cP_W)|_{V_t}\cdot P_W^{f-1},
 \qquad 0<c\le2(d-f)/\delta.
 \end{split}
\end{equation}
\end{lemma}
\begin{proof}
On the constant-rank open, the full kernels are vector bundles
with nested proper base loci. Track a component through the mark
across the $d+1$ levels. Its dimensions lie in
$\{1,\ldots,d-1\}$, so two adjacent tracked components coincide.
After a finite extension of the parameter field, spread them out
to an integral family whose marked incidence dominates $W$.
The adjacent orders satisfy
\[
 \lceil k\tau_{i+1}\rceil-\lceil k\tau_i\rceil
 \ge k\delta-1\ge h-1.
\]
Lemma~\ref{lem:jet-normal-containment} gives the ordinary-power
containment for the entire higher base ideal. Apply
Lemmas~\ref{lem:hd-base-degree} and \ref{lem:canonical-centre}
for the degree and canonical bounds. The construction follows
\cite[Lemmas~4.5--4.6]{NS}.
\end{proof}

\begin{proposition}\label{prop:hd-moving-centre}
Let $P$ be an ample $\mathbb{Q}$-Cartier divisor on $X$, and put
$\rho=(P^3)^{1/3}$. Suppose that, for some positive integer $k_*$
with $k_*P$ Cartier, a Zariski dense set of points $x$ admits
$E_x\in|k_*P|$ with $\mult_x(E_x)>4k_*\rho$.
Then there is a covering family of integral subvarieties $V$ of
one fixed dimension $d\in\{1,2\}$ such that
\[
 0<P^d\cdot V\le 3^{3-d}(P^3)^{d/3}.
\]
Moreover, on a resolution $Y\to V$, with $P_Y$ the pullback,
\begin{equation}\label{eq:moving-canonical}
 K_Y\cdot P_Y^{d-1}\le\frac{3(3-d)}{\rho}P_Y^d.
\end{equation}
The marked incidence dominates $X$.
\end{proposition}
\begin{proof}
Lemma~\ref{lem:hd-dense-multiplicity} makes the multiplicity
condition valid at every point. Taking powers of the defining
sections allows $k$ to be arbitrarily large and divisible by $k_*$.
Put
\[
 \delta=\frac{2\rho}{3},\qquad
 \tau_i=\frac32\rho+i\delta\quad(0\le i\le3),\qquad
 q_i=\lceil k\tau_i\rceil.
\]
Choose $k$ so that $k\delta\ge4$ and $k\rho/2>1$.
All full kernel spaces
\[
 V_i(x)=H^0(X,\OO_X(kP)\otimes\mathfrak m_x^{q_i})
\]
are nonzero, because $\tau_3=7\rho/2<4\rho$.
Write $B_i(x)$ for their base loci. They are proper and nested.
The point $x$ cannot be isolated in $B_0(x)$: otherwise
Lemma~\ref{lem:hd-base-degree} would give
$q_0^3\le k^3P^3$, whereas $q_0\ge3k\rho/2$.

Apply Lemma~\ref{lem:full-kernels} with $W=X$ and these four
levels. It gives a common component $V_t$ of dimension
$d\in\{1,2\}$ and a covering incidence, together with
ordinary-power containment for the entire higher base ideal,
where $h=\lceil k\delta/2\rceil$.

The higher base ideal is primary there, since $V_t$ is its component.
Lemma~\ref{lem:hd-base-degree}, applied to $H=kP$, gives
\[
 h^{3-d}k^d P^d\cdot V_t\le k^3P^3,
 \qquad
 P^d\cdot V_t\le(k/h)^{3-d}P^3
 \le(3/\rho)^{3-d}P^3.
\]
Lemma~\ref{lem:canonical-centre} gives
$K_Y\cdot P_Y^{d-1}\le(3-d)(k/h)P_Y^d$, because $K_X\sim0$.
Since $k/h\le3/\rho$, this is \eqref{eq:moving-canonical}.
Ampleness gives positivity, and the incidence is dominant.
\end{proof}

\subsubsection{The scalar bound and the movable class}\label{sec:scalar}

\begin{lemma}\label{lem:bounded-curves}
Let $Y$ be an auxiliary smooth projective variety of dimension
$f\in\{1,2,3\}$. Let $B_0$ be a reduced simple normal crossing
divisor or zero, $L$ a nef Cartier divisor of full nef dimension,
and $P_Y$ a nef $\mathbb{Q}$-Cartier divisor. If
\[
 0<P_Y^f\le C_0,\quad P_Y-rL\text{ is nef},\quad
 K_Y+B_0+L\text{ is big},
\]
and
\[
 (K_Y+B_0+L)\cdot P_Y^{f-1}\le C_1P_Y^f,
\]
then
\[
 r\le(m_fC_1)^{f-1}C_0,
\]
where $m_f$ depends only on $f$.
\end{lemma}
\begin{proof}
Birkar--Zhang \cite[Theorem~1.3]{BZ}, with boundary coefficients
in $\{0,1\}$ and Cartier index one for $L$, makes
$|m_f(K_Y+B_0+L)|$ birational. Resolve its base ideal and let
$H$ be the free moving part. Then $H$ is nef and big and
\[
 H\cdot P_Y^{f-1}\le m_fC_1P_Y^f.
\]
Khovanskii--Teissier gives
\[
 P_Y\cdot H^{f-1}\le P_Y^f
 \left(\frac{H\cdot P_Y^{f-1}}{P_Y^f}\right)^{f-1}
 \le(m_fC_1)^{f-1}C_0.
\]
General hyperplanes through the image of a very general point
under the birational morphism defined by $H$ cut an effective
curve cycle containing a curve $C$ through that point.
Full nef dimension and integrality give $L\cdot C\ge1$.
Thus $r\le P_Y\cdot C\le P_Y\cdot H^{f-1}$. For $f=1$ take $C=Y$.
This is \cite[Lemma~4.1]{NS}; in particular its constants are
independent of the varying auxiliary variety.
\end{proof}

\begin{proposition}\label{prop:scalar-bound}
Assume $\nef(X,D)=3$ and $D^3=0$. With
$P_n=r_n(D+A/n)$ as in Lemma~\ref{lem:hd-normalization}, put
$\rho_n=(P_n^3)^{1/3}$. For all sufficiently large $n$, and at a very general point $x$, every
member $E\in|kP_n|$, for every positive integer $k$ with $kP_n$
Cartier, satisfies
\[
 \mult_x(E)\le4k\rho_n.
\]
\end{proposition}
\begin{proof}
Suppose violations persist as $n\to\infty$. For fixed $n,k$, the
high-multiplicity condition is closed in the incidence space, so
a violation at a sufficiently very general point yields a
dense set of violations. Use
Proposition~\ref{prop:hd-moving-centre}, together with its canonical
bound \eqref{eq:moving-canonical}, and write $Y\to V$ for the
resolution of a general member. For $f=\dim V\in\{1,2\}$ it gives
\[
 0<P_Y^f\le3^{3-f}\rho_n^f,\qquad
 K_Y\cdot P_Y^{f-1}\le\frac{3(3-f)}{\rho_n}P_Y^f.
\]
The volume window $P_n^3\to\epsilon^3>0$ makes both bounds
uniform in $n$. Let $L$ be the pullback of $D$; then
$P_Y-r_nL$ is nef.

Choose the incidence point above the very general full-nef locus
and outside the images of all families of rational curves.
The latter exclusion is possible because $K_X\sim0$ makes $X$
non-uniruled. Since $Y\to V\subset X$ is generically finite,
$Y$ is non-uniruled and $L$ has full nef dimension; this is
the very-general-images argument of \cite[Lemma~4.2]{NS}.
By \cite{BDPP}, $K_Y$ is pseudo-effective, and
Lemma~\ref{lem:lower-dimensional} makes $K_Y+L$ big.
Also $L\cdot P_Y^{f-1}\le P_Y^f/r_n$.
Lemma~\ref{lem:bounded-curves} now bounds $r_n$ uniformly,
contrary to $r_n\to\infty$.
There are only countably many $n,k$ and jet orders, so one
very general locus handles all the stated section spaces.
\end{proof}

\begin{proposition}\label{prop:movable-from-scalar}
Assume the scalar bound at a very general point $x$, and fix
large $n$ with $\rho_n<2\epsilon$ and $r_n>s$.
Let $\beta:\widetilde X\to X$ be the blowup at $x$, with exceptional
divisor $E$. There are a birational morphism $f:V\to\widetilde X$
from a smooth projective threefold and ample Cartier divisors
$H_1,H_2$ on $V$ such that
\[
 f^*\beta^*(P+\lambda D)\cdot H_1\cdot H_2
 <40\epsilon\,f^*E\cdot H_1\cdot H_2
 \quad(0\le\lambda\le s).
\]
\end{proposition}
\begin{proof}
The class $2\beta^*P-40\epsilon E$ is not pseudo-effective.
Otherwise its convex combination with the big class $2\beta^*P$
would make $2\beta^*P-cJ$ big for rational
$8\rho_n<c<40\epsilon$. An effective divisor in a divisible
multiple of this class would push down to a member of $|2kP|$
with multiplicity exceeding the bound in
Proposition~\ref{prop:scalar-bound}.
Movable-curve duality \cite{BDPP} gives a smooth projective
birational model $f:V\to\widetilde X$ and ample real divisor classes
$H_1,H_2$ such that
\[
 2f^*\beta^*P\cdot H_1\cdot H_2
 <40\epsilon\,f^*E\cdot H_1\cdot H_2.
\]
Strict negativity permits rational approximation and clearing
denominators, so the $H_i$ may be taken ample Cartier.
For $0\le\lambda\le s$ the divisor
\[
 2P-(P+\lambda D)=(r_n-\lambda)D+r_nA/n
\]
is nef, which gives the required inequalities.
Fix this actual model and complete intersection before
reducing characteristic.
\end{proof}


\begin{proof}[Proof of Lemma~\ref{lem:full-scalar}]
Fix $s$ and take $n$ sufficiently large that
Proposition~\ref{prop:scalar-bound} applies, $\rho_n<2\epsilon$,
and $r_n>s$. Apply Proposition~\ref{prop:movable-from-scalar}
with $P=P_n$ and $\lambda=s$.
\end{proof}

\subsection{The Frobenius input}\label{sec:frobenius}

\begin{lemma}\label{lem:full-contradiction}
Let $X$ be a smooth projective threefold over $\C$ with
$K_X\sim0$ and $q(X)=0$. Let $D$ be a nef Cartier divisor with
$\nef(X,D)=3$ and $D^3=0$. Let $P$ be an ample $\mathbb{Q}$-Cartier
divisor and $s>0$ an integer such that
\[
 343P^3<\tfrac14,\qquad (2P+sD)\cdot P^2<4P^3.
\]
Let $\beta:\widetilde X\to X$ be the blowup at a point $x$, with
exceptional divisor $E$. Suppose that there are a birational morphism
$f:V\to\widetilde X$ from a smooth projective threefold and ample
Cartier divisors $H_1,H_2$ on $V$ such that
\[
 f^*\beta^*(P+sD)\cdot H_1\cdot H_2
 <\tfrac38 f^*E\cdot H_1\cdot H_2.
\]
Let $D_i,P_i$ denote pullbacks from the $i$-th factor of $X\times X$
and then to
\[
 Z=\mathbb P_{X\times X}
   \bigl(\OO_{X\times X}(D_1)\oplus\OO_{X\times X}(D_2)\bigr),
 \qquad \xi=c_1(\OO_Z(1)).
\]
Let $\pi:Z\to X\times X$ be the bundle projection, and write
$D_i=(p_i\circ\pi)^*D$, $P_i=(p_i\circ\pi)^*P$ on $Z$.
Put $R(P,s)=P_1+P_2+s\xi$, and let $U$ be the complement
of the two summand sections. For every integer $k>0$ with $kP$
Cartier, $|kR(P,s)|$ does not separate $8k$-jets at any point
of $U$.
\end{lemma}

Throughout this subsection, $X,D,P,s,x,\beta,E,f,V,H_1,H_2$ satisfy
the hypotheses of Lemma~\ref{lem:full-contradiction}. Since $f^*E$
is nonzero and effective and $H_1,H_2$ are ample,
$f^*E\cdot H_1\cdot H_2>0$. Put
\[
 b=\frac{f^*\beta^*(P+sD)\cdot H_1\cdot H_2}
 {f^*E\cdot H_1\cdot H_2}<\tfrac38.
\]
Nefness of $D$ gives, for $0\le\lambda\le s$,
\[
 f^*\beta^*(P+\lambda D)\cdot H_1\cdot H_2
 \le b\,f^*E\cdot H_1\cdot H_2.
\]
Suppose for contradiction that $|k_0R(P,s)|$ separates $8k_0$-jets
at some $z\in U$, for an integer $k_0>0$ with $k_0P$ Cartier.
All these data are fixed over $\C$ before reduction. The following
estimates specialize \cite[Theorem~6.1]{NS}; Frobenius is used
only on smooth positive-characteristic fibres.

\subsubsection{Fixed data and generic rank}

\begin{lemma}\label{lem:cotangent-flag}
For the fixed ample $\mathbb{Q}$-Cartier divisor $P$, there exist
positive integers $l_1,l_2$, with $l_iP$ Cartier, and divisors
$S\in|l_1P|$, $S_2\in|l_2P|$ such that
$X\supset S\supset C=S\cap S_2$ is a smooth integral complete
intersection flag and $\Omega_X^1|_C$ is semistable of degree zero.
In particular, $\mu_{\min}(\Omega_X^1|_C)=0$.
\end{lemma}
\begin{proof}
Choose an integer $a>0$ such that $H=aP$ is very ample Cartier.
Since $K_X\sim0$, generic cotangent semipositivity
\cite[Theorem~2.1]{CP} gives
$c_1(\mathcal Q)\cdot H^2\ge0$ for every torsion-free quotient
$\mathcal Q$ of $\Omega_X^1$. For a nonzero proper saturated
subsheaf $\mathcal F\subset\Omega_X^1$, put
$\mathcal Q=\Omega_X^1/\mathcal F$. Then
\[
 \mu_H(\mathcal F)
 =-\frac{c_1(\mathcal Q)\cdot H^2}{\operatorname{rk}\mathcal F}
 \le0=\mu_H(\Omega_X^1).
\]
Thus $\Omega_X^1$ is $H$-semistable of slope zero.

Mehta--Ramanathan restriction \cite[Theorem~6.1]{MR} provides
sufficiently large integers $m_1,m_2$ and general divisors
$S\in|m_1H|$, $S_2\in|m_2H|$ for which
$\Omega_X^1|_{S\cap S_2}$ is semistable. Bertini ensures that
$S$ and $C=S\cap S_2$ are smooth and integral. Set $l_i=am_i$.
Finally,
\[
 \deg(\Omega_X^1|_C)=K_X\cdot C=0,
 \qquad\mu_{\min}(\Omega_X^1|_C)=0.
\]
The same construction uses $\dim X-1$ general high-degree cuts
in any dimension when $K_X\equiv0$. The flag is chosen over $\C$
and fixed before the reduction in Lemma~\ref{lem:fixed-data-reduction}.
\end{proof}

\begin{lemma}\label{lem:jet-products-hd}
If sections of a line bundle $L$ on $Z$ generate jets through
order $j$ at $z$, then products of $\ell $ such sections generate jets
through order $\ell j$ for $L^\ell $.
\end{lemma}
\begin{proof}
Every monomial of degree at most $\ell j$ is a product of $\ell $ monomials
of degrees at most $j$. Lift the latter to sections with those leading
terms and multiply. The product has the desired leading term and
only higher-order errors. In a basis ordered by total degree this
gives a triangular matrix with nonzero diagonal on the finite jet
quotient, hence surjectivity.
\end{proof}

\begin{proposition}\label{prop:frob-rank-hd}
Let $Y$ be a smooth positive-characteristic fibre of a model of
$X$, and use the same notation for the specialized divisor data, over an algebraically closed field of characteristic $p$.
On $\mathbb P(\OO_Y(D)_1\oplus\OO_Y(D)_2)$, suppose a fixed
finite system of sections of $k_0(P_1+P_2+s\xi)$ generates jets
through order $8k_0$ at a point $z$ off the axes, where
$k_0P$ is Cartier. Choose $T$ so that all $T+aD$ for
$0\le a\le(k_0-1)s$ have sections nonzero at the two base points.
Put $\ell_p=\lfloor p/k_0\rfloor$, $B_p=k_0\ell_pP+T$.
For $p\ge8k_0$, the two-slot evaluation map has generic rank
$p^6$.
\end{proposition}
\begin{proof}
Lemma~\ref{lem:jet-products-hd} gives order
$8k_0\ell_p$ for $\ell_pk_0R$. Put $d_p=(p-k_0\ell_p)s$.
A weight-zero section of $d_p\xi+T_1+T_2$ nonzero at $z$ pads this
system to $pN\xi+B_{p,1}+B_{p,2}$. Its multiplication is invertible
on the local jet quotients. Moreover
\[
 8k_0\ell_p\ge8(p-k_0+1)\ge7(p-1).
\]
In seven regular parameters, a monomial surviving modulo
$(u_1^p,\ldots,u_7^p)$ has degree at most $7(p-1)$.
Thus the padded sections surject onto this Frobenius quotient.

Let $F:Y\to Y'$ be relative Frobenius, let $D'$ be the twist, and
put $\mathcal E=F_*\OO_Y(B_p)$. Smoothness makes $F$ finite flat
of degree $p^3$, so $\operatorname{rk}\mathcal E=p^3$.
The tensor product $A_*$ of the two local Frobenius-fibre algebras
has dimension $p^6$. If $t$ is the fibre coordinate and $t(z)=c\ne0$,
the full quotient is $A_*[t]/(t^p-c^p)$, free over $A_*$ with
basis $1,t,\ldots,t^{p-1}$. Projection to the coefficient of $1$
keeps exactly weights $j=pa$, with $0\le a\le s$.
Consequently the natural map
\[
 \begin{split}
 \bigoplus_{a+b=s}&H^0(Y,B_p+paD)\otimes H^0(Y,B_p+pbD)\\
 &\quad\otimes\OO_{Y'\times Y'}(-aD'_1-bD'_2)
 \longrightarrow\mathcal E\boxtimes\mathcal E
 \end{split}
\]
spans $A_*$ at the base pair. The projection formula identifies
these summands with the evaluation map; the surviving fibre scalars
$c^{pa}$ are nonzero. The map therefore has generic rank $p^6$.
No reducedness assumption on $A_*$ is used.
\end{proof}

\begin{lemma}\label{lem:fixed-data-reduction}
Fix the jet system, padding sections, canonical trivialization,
flag, blowup, and complete-intersection witness over $\C$,
with $\mu_{\min}(\Omega_X^1|_C)\ge0$. They extend to a model whose smooth closed fibres occur in
arbitrarily large characteristics and preserve the evaluation
surjection, ampleness, intersection numbers, and semistability
of $\Omega_Y^1|_C$ of slope zero.
\end{lemma}
\begin{proof}
Choose the padding divisor $T$ and the finitely many nonvanishing
padding sections as in Proposition~\ref{prop:frob-rank-hd}.
Spread out the flag, canonical trivialization, jet sections,
blowup, birational model $f:V\to\widetilde X$, and ample divisors
$H_1,H_2$. Shrinking the model preserves smoothness, ampleness,
the finite jet surjection, and all intersection numbers.
Since $\deg(\Omega_X^1|_C)=K_X\cdot C=0$, the least-slope condition
makes this curve bundle semistable of slope zero. Openness of
semistability preserves it after further shrinking.
There remain closed fibres in arbitrarily large characteristics.
\end{proof}


On each smooth fibre let $F:Y\to Y'$ denote relative Frobenius.
For the following lemmas put
\[
 \ell_p=\lfloor p/k_0\rfloor,\quad B_p=k_0\ell_pP+T,\quad
 \mathcal E_p=F_*\OO_Y(B_p),\quad s_p=p^3,\quad R_p=p^6.
\]
All constants in $O(p^{-1})$ depend only on the fixed data.

\subsubsection{The diagonal rank}

\begin{lemma}\label{lem:truncated-slope-bound}
Let $V$ be a semistable vector bundle of rank three and degree
zero on a smooth projective curve of genus $g$ in characteristic
$p$. For $0\le j\le3(p-1)$, its truncated symmetric powers satisfy
\[
 \mu_{\max}(\mathcal A_j(V))\le3c_g,
\]
where $c_g$ depends only on $g$ and the rank, independently of $j,p$.
\end{lemma}
\begin{proof}
Langer's bound
\cite[Corollary~2.5]{Langer} gives
$L_{\min}(V)\ge-c_g/(p-1)$ with $c_g$ depending only on the
genus and rank. Put $u=3(p-1)$. The perfect complementary
degree pairing gives
\[
 \mathcal A_{u-j}(V)
 \simeq\mathcal A_j(V)^\vee\otimes(\det V)^{p-1}.
\]
Truncated symmetric powers are quotients of tensor powers;
their asymptotic least slopes are at least the corresponding
multiple of $L_{\min}(V)$. Since $\deg\det V=0$, the pairing
implies $\mu_{\max}(\mathcal A_j(V))\le3c_g$, uniformly in $j,p$.
\end{proof}

\begin{lemma}\label{lem:flag-section-bound}
Let $Y\supset S\supset C$ be a smooth complete intersection flag
cut by successive ample Cartier divisors $h_1,h_2$.
If $\mathcal G$ is locally free of rank $e$ and
$\mu_{\max}(\mathcal G|_C)\le M$, with $M\ge0$, then
\[
 h^0(Y,\mathcal G)\le e(M+1)
 \prod_{i=1}^2\left(1+\frac{M}{h_i\cdot C}\right).
\]
\end{lemma}
\begin{proof}
Successive restriction along the flag bounds sections by the sum
of the spaces on $C$ twisted by $-b_1h_1-b_2h_2$, with $b_i\ge0$.
Only pairs satisfying $b_1h_1\cdot C+b_2h_2\cdot C\le M$ contribute, since
negative maximal slope forbids sections. For each pair,
evaluation at $\lfloor M\rfloor+1$ points is injective and bounds
the dimension by $e(M+1)$.
There are at most $\prod_i(1+M/(h_i\cdot C))$ pairs. The restriction
sequences terminate because sufficiently negative ample twists
have no sections.
\end{proof}

\begin{lemma}\label{lem:diagonal-rank-bound}
On a smooth fibre of Lemma~\ref{lem:fixed-data-reduction}, put
$B_p=k_0\lfloor p/k_0\rfloor P+T$ and $R_p=p^6$.
If $(2P+sD)\cdot P^2<4P^3$ and $343P^3<1/4$, the weighted evaluation
map has rank less than $R_p/4$ on the diagonal of $Y'\times Y'$.
\end{lemma}
\begin{proof}
Let $\Delta_F=Y\times_{Y'}Y$, and put
\[
 \mathcal L_p=
 \OO_{\Delta_F}(B_{p,1}+B_{p,2}+psD_1).
\]
On $\Delta_F$, the weighted section products all take values
in this same line bundle, because $pD_1=pD_2$ there.
Their value rank at any point of the diagonal $Y'\subset
Y'\times Y'$ is at most $h^0(\Delta_F,\mathcal L_p)$.
The canonical Frobenius filtration
\cite[Corollary~3.5]{KS} has graded pieces
\[
 \mathcal G_j=
 \OO_Y(2B_p+psD)\otimes\mathcal A_j(\Omega_Y^1),
 \quad0\le j\le3(p-1),\qquad
 e_j=\operatorname{rk}\mathcal A_j,\quad \sum_j e_j=p^3,
\]
Here $\mathcal A_j$ is the truncated symmetric power of degree $j$,
locally the degree-$j$ part of
\[
 \OO[u_1,u_2,u_3]/(u_1^p,u_2^p,u_3^p).
\]


By Lemma~\ref{lem:truncated-slope-bound}, the truncated powers
on $C$ have uniformly bounded maximal slopes.
Set $\Lambda=l_1l_2$, so $[C]=\Lambda P^2$ and
$h_i\cdot C=l_i\Lambda P^3$. The differences
\[
 B_p-pP=T-(p-k_0\ell_p)P
\]
belong to a fixed finite list. The displayed intersection
hypothesis therefore gives one constant $c_0$, independent of
$j,p$, for which
\[
 \mu_{\max}(\mathcal G_j|_C)
 \le M_p:=7p\Lambda P^3+c_0.
\]

Lemma~\ref{lem:flag-section-bound} therefore gives
\[
 h^0(Y,\mathcal G_j)
 \le e_jp^3\bigl(343P^3+O(p^{-1})\bigr).
\]
Summing the graded pieces gives
\begin{equation}\label{eq:diagonal-rank}
 h^0(\Delta_F,\mathcal L_p)
 \le R_p\bigl(343P^3+O(p^{-1})\bigr)<R_p/4.
\end{equation}
The last inequality follows from $343P^3<1/4$ for sufficiently
large $p$.
There is no extra factor for the number of graded pieces:
their total rank is exactly $p^3$.
\end{proof}


\subsubsection{Determinant multiplicities}

\begin{lemma}\label{lem:determinant-line-bundles}
For the evaluation map in Proposition~\ref{prop:frob-rank-hd},
put $s_p=p^3$, $R_p=p^6$, and $\mathcal E_p=F_*\OO_Y(B_p)$.
There exist line bundles
\[
 \mathcal Q_i=(\det\mathcal E_p)^{s_p}\otimes\OO_{Y'}(m_iD'),
 \qquad m_i\in\mathbb Z_{\ge0},\quad m_1+m_2=sR_p,
\]
and a nonzero determinant section of $\mathcal Q_1\boxtimes\mathcal Q_2$.
Under the numerical identification with $Y'$, for $\lambda_i=m_i/R_p$,
\[
 \frac{c_1(\mathcal Q_i)}{R_p}=P+\lambda_iD+O(p^{-1}),
 \qquad0\le\lambda_i\le s.
\]
\end{lemma}
\begin{proof}
Proposition~\ref{prop:frob-rank-hd} gives generic evaluation
rank $R_p$. Choose $R_p$ generically independent evaluation
columns of weights $(a,b)$, $a+b=s$. Their determinant is
a nonzero section
\[
 \sigma\in H^0(Y'\times Y',\mathcal Q_1\boxtimes\mathcal Q_2),
 \qquad
 \mathcal Q_i=(\det\mathcal E_p)^{s_p}\otimes\OO_{Y'}(m_iD'),
\]
where $m_i\ge0$ and $m_1+m_2=sR_p$. Let $E_\sigma$ be the
effective divisor of $\sigma$, so
\[
 E_\sigma\in|\mathcal Q_1\boxtimes\mathcal Q_2|.
\]


The Frobenius formula for the first Chern class
\cite[Lemma~4.2]{Sun} gives, under the natural numerical
identification with the twisted fibre,
\[
 \frac{c_1(\mathcal Q_i)}{R_p}
 =\frac{B_p}{p}+\frac{p-1}{2p}K_Y+\lambda_iD
 =P+\lambda_iD+O(p^{-1}),\qquad
 \lambda_i=m_i/R_p\in[0,s].
\]
\end{proof}

\begin{lemma}\label{lem:product-multiplicity-bound}
Let $L_1,L_2$ be line bundles on smooth projective varieties
$Y_1,Y_2$, and let $x_i\in Y_i$. If every $F_i\in|L_i|$ satisfies
$\mult_{x_i}(F_i)\le a_i$, then every $F\in|L_1\boxtimes L_2|$
satisfies
\[
 \mult_{(x_1,x_2)}(F)\le a_1+a_2.
\]
\end{lemma}
\begin{proof}
By K\"unneth the section space is the tensor product of the two
section spaces. Choose bases adapted to their multiplicity
filtrations. Their leading forms are independent in each degree,
and tensor products remain independent in each bidegree.
The largest possible nonzero initial degree is at most $a_1+a_2$.
\end{proof}

\begin{lemma}\label{lem:determinant-upper-bound}
Let $E_\sigma$ be the determinant divisor in
Lemma~\ref{lem:determinant-line-bundles}. The data $f,H_1,H_2$,
specialized as in Lemma~\ref{lem:fixed-data-reduction}, give
\begin{equation}\label{eq:det-upper}
 \mult_{(x',x')}(E_\sigma)
 \le R_p\bigl(2b+O(p^{-1})\bigr).
\end{equation}
\end{lemma}
\begin{proof}
Twist the specialized blowup and its fixed movable witness as well.
For every member $E_i\in|\mathcal Q_i|$, its proper transform
is effective. Pulling it back to the specialized model $V$ and
intersecting with $H_1\cdot H_2$ gives
\[
 \frac{\mult_{x'}(E_i)}{R_p}\le b+O(p^{-1}).
\]
The fixed complete intersection retains positivity on the
specialized model. The intersection inequality for effective divisors
therefore applies to every member of these linear systems on the fibre.

Apply Lemma~\ref{lem:product-multiplicity-bound} to the two
line bundles $\mathcal Q_i$.
\end{proof}

\begin{lemma}\label{lem:matrix-multiplicity-bound}
Let $M$ be an $R\times R$ matrix of regular functions near a
smooth point $u$, with $\det M\ne0$ and $\operatorname{rank}M(u)=r$.
Then
\[
 \mult_u(\operatorname{div}(\det M))\ge R-r.
\]
\end{lemma}
\begin{proof}
Constant invertible row operations make $R-r$ rows vanish at $u$.
Each determinant monomial contains one entry from each such row,
so belongs to $\mathfrak m_u^{R-r}$.
\end{proof}



\begin{proof}[Proof of Lemma~\ref{lem:full-contradiction}]
Choose the flag from Lemma~\ref{lem:cotangent-flag}.
Lemma~\ref{lem:fixed-data-reduction} spreads the fixed flag,
jet system, and curve class to smooth fibres in arbitrarily
large characteristics $p$. Proposition~\ref{prop:frob-rank-hd}
and Lemma~\ref{lem:determinant-line-bundles} give a nonzero
determinant divisor $E_\sigma$, with $R_p=p^6$.
The diagonal rank bound in Lemma~\ref{lem:diagonal-rank-bound},
the matrix bound in Lemma~\ref{lem:matrix-multiplicity-bound},
and Lemma~\ref{lem:determinant-upper-bound} give
\begin{equation}\label{eq:det-lower}
 \tfrac34R_p<\mult_{(x',x')}(E_\sigma)
 \le R_p\bigl(2b+O(p^{-1})\bigr).
\end{equation}
Divide by $R_p$ and let $p\to\infty$. Then
$3/4\le2b<3/4$, a contradiction.
\end{proof}

\begin{proof}[Proof of Lemma~\ref{lem:full-nojets}]
Fix $s>0$ and assume $\kappa(D)=-\infty$. Take $n$ sufficiently large
for Lemma~\ref{lem:full-scalar}, $r_n>s$ and $343P_n^3<1/4$.
Since $P_n-r_nD$ is nef,
\[
 (2P_n+sD)\cdot P_n^2\le(2+s/r_n)P_n^3<4P_n^3.
\]
Lemma~\ref{lem:full-scalar} supplies a point $x$, its blowup
$\beta:\widetilde X\to X$ with exceptional divisor $E$, and
$f:V\to\widetilde X$, $H_1,H_2$ such that
\[
 f^*\beta^*(P_n+sD)\cdot H_1\cdot H_2
 <40\epsilon f^*E\cdot H_1\cdot H_2
 <\tfrac38 f^*E\cdot H_1\cdot H_2.
\]
Apply Lemma~\ref{lem:full-contradiction} with $P=P_n$.
Since $R(P_n,s)=R_{n,s}$, no $|kR_{n,s}|$ with $kR_{n,s}$ Cartier
separates $8k$-jets at any $z\in U$.
\end{proof}

\begin{thebibliography}{99}
\bibitem{LA}
\emph{Log abundance in characteristic zero},
OpenAI mathematics collection, September 24, 2026,
Lemmas~7.1--7.5 (Tools for moving base components).
\href{https://github.com/openai/math/blob/main/preprints/Log-abundance-in-characteristic-zero-September-24-2026/paper.pdf}{Source manuscript}.
\bibitem{FH}
O.~Fujino and K.~Hashizume,
\emph{Existence of log canonical modifications and its applications},
Eur. J. Math. \textbf{9} (2023), Paper No.~13, Theorem~2.10.
\href{https://doi.org/10.1007/s40879-023-00598-0}{DOI}.
\bibitem{NS}
\emph{Numerical Semiampleness of Nef Adjoint Divisors},
OpenAI mathematics collection, October 3, 2026,
Proposition~3.3, Lemmas~4.1--4.7, Proposition~5.3,
Theorem~6.1, and Sections~8.4--8.5.
\href{https://github.com/openai/math/blob/main/preprints/Numerical-Semiampleness-of-Nef-Adjoint-Divisors-October-3-2026/paper.pdf}{Source manuscript}.
\bibitem{LP}
V.~Lazi\'c and Th.~Peternell,
\emph{On Generalised Abundance, I},
Publ. Res. Inst. Math. Sci. \textbf{56} (2020), 353--389.
Numbering follows arXiv:1808.00438v2, especially Corollary~C
and Lemmas~2.14 and 3.1.
\href{https://arxiv.org/pdf/1808.00438}{Preprint}.
\bibitem{FG}
O.~Fujino and Y.~Gongyo,
\emph{On the moduli $b$-divisors of lc-trivial fibrations},
Ann. Inst. Fourier \textbf{64} (2014), 1721--1735,
Definition~3.1 and Theorem~3.5.
\href{https://doi.org/10.5802/aif.2894}{DOI}.
\bibitem{BZ}
C.~Birkar and D.-Q.~Zhang,
\emph{Effectivity of Iitaka fibrations and pluricanonical
systems of polarized pairs},
Publ. Math. Inst. Hautes \'Etudes Sci. \textbf{123} (2016),
283--331, Theorem~1.3.
\href{https://doi.org/10.1007/s10240-016-0080-x}{DOI}.
\bibitem{BDPP}
S.~Boucksom, J.-P.~Demailly, M.~P\u{a}un, and Th.~Peternell,
\emph{The pseudo-effective cone of a compact K\"ahler manifold and
varieties of negative Kodaira dimension},
J. Algebraic Geom. \textbf{22} (2013), 201--248.
\href{https://doi.org/10.1090/S1056-3911-2012-00574-8}{DOI}.
\bibitem{CP}
F.~Campana and M.~P\u{a}un,
\emph{Orbifold generic semi-positivity: an application to families
of canonically polarized manifolds},
Ann. Inst. Fourier \textbf{65} (2015), 835--861, Theorem~2.1.
\href{https://doi.org/10.5802/aif.2945}{DOI}.
\bibitem{MR}
V.~B.~Mehta and A.~Ramanathan,
\emph{Semistable sheaves on projective varieties and their
restriction to curves},
Math. Ann. \textbf{258} (1982), 213--224, Theorem~6.1.
\href{https://doi.org/10.1007/BF01450677}{DOI}.
\bibitem{Langer}
A.~Langer,
\emph{Semistable sheaves in positive characteristic},
Ann. of Math. \textbf{159} (2004), 251--276, Corollary~2.5.
\href{https://doi.org/10.4007/annals.2004.159.251}{DOI}.
\bibitem{KS}
Y.~Kitadai and H.~Sumihiro,
\emph{Canonical filtrations and stability of direct images by
Frobenius morphisms},
Tohoku Math. J. \textbf{60} (2008), 287--301, Corollary~3.5.
\href{https://doi.org/10.2748/tmj/1215442876}{DOI}.
\bibitem{Sun}
X.~Sun,
\emph{Direct images of bundles under Frobenius morphism},
Invent. Math. \textbf{173} (2008), 427--447,
Theorem~3.7 and Lemma~4.2; numbering follows
arXiv:math/0611360v2.
\href{https://doi.org/10.1007/s00222-008-0125-y}{DOI}.
\end{thebibliography}

\end{document}
